\documentclass[10pt]{article}

\usepackage[utf8]{inputenc}

\usepackage[T1]{fontenc}

\usepackage{amsmath}

\usepackage{amsfonts}

\usepackage{amssymb}

\usepackage[version=4]{mhchem}

\usepackage{stmaryrd}

\usepackage{bbold}

\usepackage{graphicx}

\usepackage[export]{adjustbox}

\graphicspath{ {./images/} }



\begin{document}

Тогда $A_{n}=A$ и матрица $A^{-1}$ получается в результате $n$-кратного применения описанного выше процесса. Расчетные формулы при этом имеют следуюпий вид: если $a_{j}^{(k)}$ есть $j$-й столбец $A_{k}$, то для $k=1,2, \ldots, n$;



\[

a_{j}^{(k)}=a_{j}^{(k-1)}-\frac{a_{k} a_{j}^{(k-1)}}{1+a_{k} a_{k}^{(k-1)}} a_{k}^{(k-1)}, j=1,2, \ldots, n .

\]



Для матрицы $A_{k}^{-1}$ достаточно вычислять элементы первых $k$ строк, т. к. последующие строки совпадают со строками единичной матрицы.



Известны другие способы организации вычислений в П. м., основанные на модификации (*), напр. т. н. м е т о д E p III о в а (cM. [1]).



Лum.: [1] $\Phi$ а д д е е в Д. К., $\Phi$ а д д е е в а В. Н., Вычислительные методы линейной алгебры, 2 изд., М. $\frac{1}{\Gamma}$. Д., 1963 . Ки.



ПОРИСТОСТИ ТОЧКА для множества $E$ из $n$-мерного евклидова пространства $\mathbb{R}^{n}-$ точка $x_{0} \in \mathbb{R}^{n}$, $n \geqslant 1$, для к-рой существует последовательность открытых шаров $B_{k}$ с радиусами $r_{k} \rightarrow 0$ и общим центром в точке $x_{0}$ таких, что для каждого $k=1,2, \ldots$ найдется открытый шар $G_{k} \subset B_{k} \backslash E$ радиуса $\geqslant C r_{k}$, где $C$ положительно и не зависит от $k$ (но, вообще говоря, зависит от $x_{0}$ и $E$ ). Множество $E$ наз. п о р и с т ы м, если каждая его точка является П. т. для него. Множество $E$ наз. $\sigma-\Pi$ о р и с т ы м, если его можно представить в виде конечного или счетного объединения пористых множеств (см. [1]). П. т. для $E$ является П. т. для его замыкания $\bar{E}$ и не является точкой плотвости в смысле Лебега ни для $E$, ни для $\bar{E}$. Каждое пористое или $\sigma$-пористое множество $E \subset \mathbb{R}^{n} \quad$ имеет первую категорию по Бэру и нулевую меру Лебега в $\mathbb{R} n$. Обратное, вообще говоря, неверно: существуют даже совершенные нигде не плотные множества $E \subset \mathbb{R} \mathbf{1}$, имеющие меру нуль, но не являющиеся $\sigma$-пористыми (см. [2]). Для множества $E$, лежащего на гладком многообразии $S \subset \mathbb{R}^{n}, \Pi$. т. $x_{0} \in S$ множества $E$ относительно многообразия $S$ определяется, как выше, при дополнительном условии, что центры шаров $B_{k}$ лежат на $S$.



\includegraphics[max width=\textwidth, center]{2023_07_08_c8d3b4eb56693f20d3e0g-1}

pis pěst. mat.», 1976, sv. 101, s. 350-59. E. П. Долженко.



ПОРЦИЯ м н о ж е т в а - пересечение множества с интервалом в случае множества на прямой и с открытым кругом или шаром, с открытым прямоугольником или параллелепипедом в случае множества в $n$-мерном $(n \geqslant 2)$ пространстве. Важность этого понятия оправдывается следующими обстоятельствами. Множество $A$ является всюду плотным в множестве $B$, если каждая П. множества $B$ содержит точку множества $A$, иначе говоря, если замыкание $\bar{A} \supset B$. Множество $A$ является нигде не плотным в множестве $B$, если $A$ не будет всюду плотным ни в какой П. множества $B$, т. е. если не существует П. множества $B$, содержащейся в $\bar{A}$. А. А. Конюшиов.

с р а в н е н и е ПОРЯДКА СООТНОІІЕНИЕ, с р а в н е н и е фу укци й, $O-0$-соотношен ия, ас импто тиче с к и е с о о тношен и я, - понятие, возникающее при изучении поведения одной функции относительно другой в окрестности нек-рой точки (быть может, бесконечной).



Пусть $x_{0}$ - предельная точка множества $E$. Если для функций $f(x)$ и $\boldsymbol{g}(x)$ существуют такие постоянные $c>0$ и $\delta>0$, что $|f(x)| \leqslant c|g(x)|$ при $\left|x-x_{0}\right|<\delta$, $x \neq x_{0}$, то говорят, что $f$ является ограниченной по сравнению c $g$ функцией в нек-рой окрестности точки $x_{0}$, и пишут



\[

f(x)=O(g(x)) \text { при } x \longrightarrow x_{0}

\]



(читается: "f $(x)$ есть $O$ больное от $g(x) ») ; x \rightarrow x_{0}$ означает, что рассматриваемое своїство имеет место лишь в нек-рой окрестности точки $x_{0}$. Естественным образом это определение переносится на случай $x \rightarrow \infty, x \rightarrow$ $\rightarrow \pm \infty$.



Если функции $f(x)$ и $g(x)$ такие, что $f=O(g)$ и $g=O(f)$ при $x \rightarrow x_{0}$, то они наз. фу н к гу ия м и о дн о г о п о р я д к а прІі $x \rightarrow x_{0}$. Напр., если функции $\alpha(x)$, $\beta(x)$ таковы, что $\alpha(x) \neq 0, \beta(x) \neq 0$ при $x \neq x_{0}$ и существует предел



\[

\lim _{x \rightarrow x_{0}} \frac{\alpha(x)}{\beta(x)}=c \neq 0

\]



то они одного порядка при $x \rightarrow x_{0}$.



Две функции $f(x)$ и $g(x)$ наз. э к в и в а л е н т ным И (асимптотически равными) при $x \rightarrow x_{0}$ и пишут $f \sim g$, если в нек-рой окрестности точки $x_{0}$, кроме, быть может, самой точки $x_{0}$, определена такая функция $\varphi(x)$, что



\[

f(x)=\varphi(x) g(x) \text { и } \lim _{x \rightarrow x_{0}} \varphi(x)=1 .

\]



Условие әквивалентности двух функций симметрично, т. е. если $f \sim g$, то и $g \sim f$ при $x \rightarrow x_{0}$, и транзитивно, т. е. если $f \sim g$ и $g \sim h$, то $f \sim h$ при $x \rightarrow x_{0}$. Если в нек-рой окрестности точки $x_{0}$ при $x \neq x_{0}$ справедливы неравенства $f(x) \neq 0, \quad g(x) \neq 0$, то условие (\%) эквивалентно любому из следующих



\[

\lim _{x \rightarrow x_{0}} \frac{f(x)}{g(x)}=1, \lim _{x \rightarrow x_{0}} \frac{g(x)}{f(x)}=1 .

\]



Если $\alpha(x)=\varepsilon(x) f(x)$, где $\lim _{x \rightarrow x_{0}} \varepsilon(x)=0$, то говорят, что $\alpha$ является б е с конечн о ма ло й ф у нк ци е й по сравнению с функцией $f$, и пишут



\[

\alpha=o(f) \text { при } x \rightarrow x_{0}

\]



(читается: « $\alpha$ есть $о$ малое от $f »)$. Если $f(x) \neq 0$ при $x \neq x_{0}$, то $\alpha=o(f)$, когда $\lim _{x \rightarrow x_{0}} \alpha / f=0$. В случае, если $f$ является бесконечно малой при $x \rightarrow x_{0}$, то говорят, что функция $\alpha=o(f)$ при $x \rightarrow x_{0}$ является б е с коне ч но м ало й боле е вы с о к о г о по р я д ка, чем $f$. Если же $g(x)$ и $[f(x)]^{k}-$ величины одного порядка, то говорят, что $g$ является в елич И но й п о р я дк а $k$ относительно $f$. Все формулы указанного вышІе вида наз. а с и м п тотиче с кими о це н к ам и, они наиболее интересны для бесконечно малых и бесконечно больших функций.



П р и м е р ы: $\quad e^{x-1}=o(1)(x \rightarrow 0) ; \quad \cos x^{2}=O(1)$, $(\ln x)^{\alpha=o\left(x^{\beta}\right)}(x \rightarrow \infty ; \alpha, \beta-$ любые положительные числа); $[x / \sin (1 / x)]=O\left(x^{2}\right)(x \rightarrow \infty)$.



Некоторые свойства символов $o$ и $O$ :



\[

\begin{gathered}

O(\alpha f)=O(f) \quad(\alpha-\text { const }) \\

O(O(f))=O(f) \\

O(f) O(g)=O(f \cdot g) \\

O(o(f))=o(O(f))=o(f) \\

O(f) o(g)=o(f \cdot g)

\end{gathered}

\]



если $0<x<x_{0}$ и $f=O(g)$, то



\[

\int_{x_{0}}^{x} f(y) d y=O\left(\int_{x_{0}}^{x}|g(y)| d y\right)\left(x \longrightarrow x_{0}\right) .

\]



Формулы, содержащие $o-, O$-символы, читаются только слева направо, это не исключает того, что отдельные формулы оказываются справедливыми и при чтении справа налево. Символы о и $O$ для функций комплексного переменного и для функций многих переменных вводятся аналогично тому, как они были определены выше для функций одного действительного переменного.



ПОРЯДКОВАЯ СТАТИСТИКА - член вариационного ряда, построенного по результатам наблюдений. Пусть





\end{document}